% Generated draft: requires mathematical and source review.
\documentclass[11pt]{article}
\usepackage[margin=25mm]{geometry}
\usepackage{fontspec}
\setmainfont{Latin Modern Roman}
\usepackage{amsmath,amssymb,mathtools,unicode-math}
\setmathfont{Latin Modern Math}
\usepackage{enumitem,needspace,xurl,hyperref}
\setlength{\parindent}{0pt}
\setlength{\parskip}{5pt}
\setlength{\emergencystretch}{3em}
\newcommand{\N}{\mathbb{N}}
\newcommand{\Z}{\mathbb{Z}}
\newcommand{\Q}{\mathbb{Q}}
\newcommand{\R}{\mathbb{R}}
\newcommand{\C}{\mathbb{C}}
\newcommand{\F}{\mathbb{F}}
\newcommand{\A}{\mathbb{A}}
\newcommand{\PP}{\mathbb{P}}
\newcommand{\E}{\mathbb{E}}
\newcommand{\eps}{\varepsilon}
\newcommand{\dd}{\,\mathrm{d}}
\begin{document}
\section*{UnsolvedMath ID: 30005437}
\textbf{Source problem code:} catalogue entry
\subsection{Definitions and mathematical statement}
The supplied record does not define a \(qL\)-algebra, its monoid-valued or group-valued invariants, or the intended meaning of a Garside-like structure. These terms therefore remain unresolved parts of the proposed research programme; no axioms or universal construction can be recovered from this record alone.

The supplied motivation states that non-degenerate set-theoretic solutions of the Yang--Baxter equation are parametrized by non-degenerate \(q\)-cycle sets, while non-degenerate involutive solutions are parametrized by non-degenerate cycle sets. It further attributes to Rump an association of an \(L\)-algebra to each non-degenerate cycle set, leading to an approach to Garside structures on the structure monoids and groups of cycle sets. Definitions of these objects and the details of that association are not supplied.

The research question, identified as Problem 11 in the supplied excerpt, is:
\begin{quote}
Develop a theory of \(qL\)-algebras and their monoid-valued and group-valued invariants. Determine whether these invariants carry Garside-like structures.
\end{quote}
The admissible class of \(qL\)-algebras, the constructions of the invariants, the required Garside-like properties, and any hypotheses under which those properties are sought are unspecified. In particular, the record does not determine whether the final question concerns every object of an intended class or a restricted subclass. This is consequently a research-programme statement with unresolved definitions, rather than a fully quantified conjecture.
\subsection{Short English statement}
Develop the proposed theory of \(qL\)-algebras, construct its monoid-valued and group-valued invariants, and investigate whether they have Garside-like structures. The supplied record leaves the defining axioms, invariant constructions, and intended structural properties unspecified.
\subsection{Sources}
\begin{itemize}
\item Local UnsolvedMath registry: \texttt{data/unsolvedmath/problems.json}, permanent ID \(30005437\), title ``Monoid and Group Invariants of \(qL\)-Algebras,'' exact source code \texttt{catalogue entry}. The supplied record assigns the problem to the Oberwolfach Reports Open Problems collection and gives proposed year 2023. No individual proposer is supplied.
\item Supplied report citation: \emph{Mini-Workshop: Skew Braces and the Yang--Baxter Equation} (2023), \emph{Oberwolfach Reports}, pp.\ 537--563, DOI \texttt{10.4171/owr/2023/9}, \url{https://doi.org/10.4171/owr/2023/9}. The supplied excerpt labels the target as Problem 11. The literature-source entry labelled ``original OWR report'' supplies the same link: \url{https://doi.org/10.4171/owr/2023/9}.
\item Contextual attribution: the supplied excerpt credits Rump with associating an \(L\)-algebra to each non-degenerate cycle set. It also refers to reference [8] for the parametrization by non-degenerate \(q\)-cycle sets; the bibliographic identity of [8] is not supplied.
\item Supplied dated literature curation, checked \texttt{2026-08-22}: status ``open.'' Its assessment concerns whether displayed relations admit a brace-like theory, states that those relations are absent from the extraction, and reports that no verified general construction or classification was found. This is supplied dated curation, not an independently verified literature assessment; its relationship to the \(qL\)-algebra target is unclear. The background explicitly requests independent verification of time-sensitive status and cited claims.
\item The supplied registry labels the statement \texttt{corrected\_verified} and reports that the official report was checked directly and that the clean statement restores definitions and removes extraction artifacts. These are supplied verification claims, not verification performed for this draft; the supplied clean statement contains no defining axioms for \(qL\)-algebras.
\end{itemize}
\end{document}
